% =======================================
% MEETING 10
% ---------------------------------------
% DATE   - September 23, 2026
% FORMAT - F2F
% =======================================
% TOPIC  - Aspects of Special Relativity
% =======================================


\subsection{Elements of Special Relativity}
	A lot of the mathematics done in General Relativity take 
	inspiration from Special Relativity.
	In fact, Special Relativity is a special case of GR.

	Everything we did in \S \ref{sec:Aspects-of-Differential-Geometry}
	is to describe the mathematical foundations of GR, so we can 
	also discuss SR using the tools we learned. We could have learned
	SR first, just as we do in Introductory Physics II, using the 
	elementary assumptions and definitions, and arrive with the 
	concepts that arise from there, then use those tools to understand
	GR\footnote{In fact, that's what they did originally, Einstein and the others.}
	However, we can more aptly describe SR using the tools of differential geometry,
	which is why we spent so much time motivating ourselves to learn it.

	\subsubsection{$n$-Dimensional Euclidean Space}
		We start, as we always do, with the familiar notions and preconceptions
		we have about the mathematical nature of things, but with a little more specificity.

		We describe space intuitively as a three-dimensional Euclidean manifold
		$\R^3$, since we are creatures of three spatial dimensions. We can
		generalize this notion of Euclidean space to $n$ dimensions.

		For a manifold $\mcal M = \R^n$, we equip it with the metric pairing
		\[
			(\R^n,\delta),
		\]
		where $\delta$ is the Euclidean metric. In terms of the coordinate basis,
		the metric is
		\begin{equation}
			\label{eq:euclidean-metric}
			\delta
			=
			\delta_{\mu\nu}
			\dd{x^\mu}
			\otimes
			\dd{x^\nu}.
		\end{equation}
		The components of the Euclidean metric are
		\[
			\delta_{\mu\nu}
			=
			\diag(1,\cdots,1),
		\]
		where $1$ is repeated $n$ times.

		We are not restricted to any bizarre coordinate system. In fact,
		the Euclidean metric is the natural metric on $\R^n$, since it
		corresponds to how we intuitively conceive of ordinary flat space.

		We can naturally use the Cartesian coordinates
		$(x,y,z,\cdots,x^n)$ to parameterize the manifold,
		\[
			\psi
			=
			(x,y,z,\ldots,x^n),
			\qquad
			x,y,z,\ldots,x^n\in\R.
		\]
		Installing the Euclidean metric on the manifold gives rise to many
		of the familiar concepts from elementary vector analysis and real
		analysis, such as the vector
		\[
			\vb{v}
			=
			\bmqty{
				v^1\\
				v^2\\
				\vdots\\
				v^n
			},
		\]
		the dot product,
		\[
			\vb{u}\cdot\vb{v}
			=
			\delta_{\mu\nu}u^\mu v^\nu,
		\]
		the cross product in three dimensions,
		\[
			(\vb{u}\times\vb{v})^\sigma
			=
			\varepsilon\indices{_{\mu\nu}^\sigma}
			u^\mu v^\nu,
		\]
		and the other concepts which stem from this structure.

		For all of this, however, time is considered a separate entity,
		assumed to be absolute, with the measurement of time being the same
		for all observers.

		However, as we learned more about the nature of space and time,
		we realized that this is not the fundamental description of reality.
		There is a more natural way to describe reality which unifies both
		space and time. This is the Minkowski metric.

	\subsubsection{Minkowski Spacetime}

		The Minkoswki metric is a Semi-Riemannian metric, specifically a 
		Lorentzian metric, i.e., only one diagonal component is negative.

		The Minkowski metric is defined for $n$-dimensional manifold.
		For describing the manifold $\mcal M= \R^n$, we equip it with tne 
		metric pairing
		\[
			(\R^n,\eta)
		\]
		Where $\eta$ is defined to be
		\[
			\eta \equiv \eta_{\mu\nu}\dd{x^\mu}\otimes \dd{x^\nu}
		\]
		and the components are
		\[
			\eta_{\mu\nu} = \diag(-1,1,\cdots,1)
		\]
		where the $+1$ components are repeated $n-1$ times.

		For GR, we normally restrict ourselves to the 4D Minkowski metric.
		Hence, we define:
		\begin{definition}{Minkowski Metric}
			The Minkowski metric $\eta$ is the metric tensor of flat spacetime
			whose components in an inertial coordinate system are
			\[
				\eta_{\mu\nu}
				=
				\diag(-1,1,1,1),
			\]
		\end{definition}
		
		Where we can naturally install the Lorentzian coordinates
		\[
			\qty{
				x^\mu
			}
			=
			(ct,x,y,z)
		\]
		This is the Minkowski spacetime, it is defined with the Minkowski metric.

		The entirety of special relativity can be described as nothing 
		more than physics on Minkowski spacetime.

		For general relativity, we are no longer restricted on the Minkowski metric.
		However, with the principle of equivalence, just like how any Riemannian manifold is locally flat, or locally Euclidean, any Lorentzian manifold, if you zoom in enough,
		is also locally flat, i.e., it locally approximates the Minkowski metric.
		
		So physics in general relativity is locally special relativity, which
		is why a lot of the concepts and aspects learned from SR generalizes naturally to GR.
		\newpage

	\subsubsection{Concepts and Notions}
		
		Special Relativity is built upon several fundamental concepts that
		we must first define before we can discuss how space and time are
		measured. In particular, we need to specify what we mean by an
		observer, a particle, a world line, and a reference frame. These
		concepts describe how physical systems move through spacetime and
		how they assign coordinates and measurements to physical events.

		\begin{definition}{Event}
			An event is defined to be a point in
			space and time, $(t,x,y,z)$. A collection
			of all the events is called a \textbf{spacetime},
			and thus each event is a spacetime point.
		\end{definition}

		An event is, in general, just a point in a 4-dimensional
		manifold (with proper additional structures such as a Lorentzian metric).
		An example of an event is a supernovae. It happened at a particular
		time $t_\star$, and at a position $(x_\star,y_\star, z_\star)$.
		Therefore, this particular event can be descrived as a point, $x^\mu$,
		such that
		\[
			x^\mu = (t_\star, x_\star, y_\star, z_\star)
		\]
			
		\begin{definition}{Observer}
			An observer is an idealized physical system which carries an
			instrument to measure its local passing of time, i.e, a clock
			and measuring apparatus, and which makes observations and
			measurements along its world line.
		\end{definition}
		To make a measurement, the observer should
		be equipped with an accurate clock, called a \textbf{standard clock}, and the reading of this
		clock is called the \textbf{proper time} of this observer (see \S\ref{subsubsec:proper-time} for details). More
		generally, one can consider that any point mass carries a standard clock, and each
		point mass has its own proper time. Additionally, an observer must 
		have an instrument which allows it to measure distances and displacements relative
		to itself.

		An observer can only make direct measurements at its own location in
		spacetime. In other words, all measurements made by an observer are
		local; the observer can directly measure only quantities associated
		with the event in its own world line.
		It cannot measure events on other locations instantaneously. 

		An observer can nevertheless obtain information about events that
		occur elsewhere. This is possible only when some physical
		interaction carrying information from that event reaches the
		observer. For example, an observer may receive a photon emitted by
		a distant star, or a gravitational wave produced by a distant
		astrophysical event. The observer does not directly observe the
		distant event itself; rather, it observes the information carried to
		it by the physical interaction long after the event actually happened.

		As we have shown in \S\ref{subsec:relativity-without-light}, there is an invariant speed $c$ in spacetime,
		which represents the maximum speed at which causal influences and
		information can propagate through spacetime. Light travels at this
		speed in vacuum, and any physical signal carrying information must
		propagate at a speed no greater than $c$.

		Consequently, the events which can have a causal relationship with
		a given observer are restricted by the finite speed of propagation.
		This restriction is represented geometrically by the observer's
		light cone.

		\begin{definition}{Light cone}
			The light cone of an event is the set of spacetime events that
			can be connected to that event by a light signal traveling at
			the invariant speed $c$. (See fig. \ref{fig:lightcone})

			\hspace{1em}The future light cone consists of events that can be causally
			influenced by a signal emitted from the event, while the past
			light cone consists of events from which a light signal could
			have reached the event.
		\end{definition}

		For example, consider an observer located at the spatial origin
		at $t=0$. A light signal emitted by the observer can travel a
		distance $r=ct$ after a time $t$. Therefore, in a spacetime diagram,
		the boundary of the future light cone is described by%
		\[
			r=ct.
		\]
		Similarly, a light signal arriving at the observer at $t=0$ must
		have been emitted at a spatial distance $r=-ct$ at the corresponding
		past time. These boundaries separate spacetime into regions with
		different causal relationships to the event.%
		\footnote{A light-cone is really a 4D structure. However, in its 
		3D projection, we visualize it as a cone, hence the name.}
\newpage
		\begin{figure}[h!]
			\centering
			\includegraphics[width=3in, height=3in]{images/lightcone.pdf}
			\caption{A light cone showing the causal structure of spacetime.
			The $z$ axis is contracted. Events on $-t$ represent the past,
			and events on $+t$ represent the future. It is the causal boundary
			of an observer.}
			\label{fig:lightcone}
		\end{figure}
	
		An observer can receive information only from events lying within
		or on its past light cone. Conversely, an observer can causally
		influence only events lying within or on its future light cone.
		Events outside these light cones are said to be spacelike separated
		from the observer's event. No signal traveling at or below $c$ can
		carry information between such events.

		Thus, the light cone expresses a fundamental causal structure of
		Special Relativity. It determines which events can influence one
		another and which events cannot have a causal connection.
		We can then reframe $c$ as the \textit{speed of causality}.


		\begin{definition}{Particle}
			A particle is an idealized physical point mass object whose spatial extent
			is neglected, so that it may be represented as a point in space
			at each instant of time. Its motion through spacetime is therefore
			represented by a world line.
		\end{definition}

		Particles may be classified according to their invariant mass.
		They can be either massive or massless:
		\begin{itemize}
			\item \text{Massive particles}

				A massive particle possesses a nonzero invariant mass.
				Examples include electrons, atoms, spacecraft, planets,
				stars, and you. In the idealized point-particle approximation,
				the spatial extent of the object is neglected and its
				motion is represented by a single world line.

			\item \text{Massless particles}

				A massless particle has zero invariant mass. The photon
				is the familiar example. In Special Relativity, a
				massless particle propagates at the invariant speed $c$
				in vacuum and follows a null trajectory through
				spacetime.
		\end{itemize}

		An observer need not be identified with a fundamental particle.
		An observer can be any physical system capable of carrying a clock
		and performing measurements. For example, a spacecraft, laboratory,
		or person can be treated as an observer. When an observer is
		idealized as a pointlike object, its motion through spacetime can
		be represented by a timelike world line.

		\begin{definition}{World line}
			The world line of a particle or observer is the curve in
			spacetime that represents its entire history.

			\hspace{1em}At each instant of its proper time, the particle occupies a
			particular event in spacetime. As the particle moves through
			spacetime, these events form a continuous curve, called its
			world line.
		\end{definition}

		The world line is therefore the spacetime analogue of the
		trajectory of a particle through ordinary space. In three-dimensional
		space, we may describe a particle's trajectory by
		$\vb{r}(t)$. In spacetime, however, we describe its history by a
		curve
		\[
			x^\mu(\tau),
		\]
		where $\tau$ is the particle's proper time for a massive particle,
		The parameter $\tau$ measures the time recorded by a clock moving
		along the particle's world line, and each point $x^\mu{\tau}$ is
		an event.

		For an inertial observer, the world line is a straight timelike
		line in Minkowski spacetime. Equivalently, it is a geodesic of
		Minkowski spacetime. An observer that accelerates, on the other
		hand, follows a non-geodesic world line.

		Representing a world line in a graph is called a \textbf{spacetime diagram}.
		In a spacetime diagram,
		the upward direction represents the time direction, and the horizontal directions
		represent spatial directions. Each horizontal slice represents the whole space at a
		certain moment of time. One can see the whole process of the motion (evolution) by
		viewing the spacetime diagram from the bottom up.

		\begin{figure}[h]
			\centering
			\includegraphics[width=4in,height=3in]{images/wordline.pdf}
			\caption{The trajectory of a particle in $\R^3$ and its
			corresponding world line in spacetime. This is a spacetime diagram. The spatial coordinate
			$z$ has been contracted for visualization.}
			\label{fig:worldline}
		\end{figure}

		\begin{definition}{Reference frame}
			A reference frame is a system of coordinates and clocks used by
			an observer to assign spatial and temporal coordinates to events.
			It provides a particular viewpoint from which measurements of
			spacetime events can be described.
		\end{definition}

		A reference frame allows an observer to assign coordinates
		$(t,x,y,z)$ to events. Different observers may use different
		reference frames and therefore assign different spatial and
		temporal coordinates to the same event. Special Relativity
		describes precisely how these coordinates are related when
		observers move relative to one another.

		An inertial reference frame is a reference frame associated with
		an inertial observer, or more generally with a family of observers
		who are all in uniform relative motion and whose coordinate
		system is non-accelerating. In Minkowski spacetime, the spatial
		coordinate axes of such a frame may be regarded as being carried
		by an idealized congruence of observers whose world lines are
		straight and parallel.
		All inertial observers are on an equal footing, i.e., no inertial observer is pre-
		ferred over any other; that is, one cannot choose a special element (or several) from
		the subset formed by inertial observers. For example, one cannot ask which inertial
		observer is at absolute rest.	

		A reference frame is defined abstractly as an infinite set 
		of observers (each carrying a standard clock and labeled by spatial 
		coordinates) such that every point in spacetime — or in an open subset
		of it — is passed through by exactly one observer, meaning each observer's 
		worldline is unique and non-intersecting. This formalizes and generalizes
		the intuitive notion of a frame. For example, the familiar ``train frame" imagines 
		observers filling all of space and moving together with the train, while 
		the ``ground frame" is another such set with a relative velocity, and in 
		a spacetime diagram their worldlines appear as parallel oblique 
		lines versus vertical lines. The definition also permits more general 
		frames where worldlines need not be parallel, allowing distances between 
		observers to change over time.

		Thus, we may think of a reference frame as an extended network
		of idealized observers, each carrying a clock and occupying a
		different spatial position. Together, these observers provide a
		way of assigning coordinates to events throughout spacetime.
		The observer at a particular spatial position measures its own
		local proper time, while the collection of synchronized clocks
		allows the reference frame to assign a common coordinate time to
		events occurring at different locations.

		\newpage

	\subsubsection{Lorentz Transformations}
		\label{subsubsec:lorentz-transformations}

		Suppose two observers in spacetime, with curves (or worldlines)
		\[
			\qty{x'^\mu}\qquad \qty{x^\nu}
		\]

		Transformations between the two frames of reference are:
		\begin{equation}
			\label{eq:lorentz-transformation-idx}
			x'^\mu = \Lambda\indices{^\mu_\nu}x^\nu
		\end{equation}
		and
		\begin{equation}
			\label{eq:lorentz-inverse-transform}
			x^\nu  = (\Lambda\inv)\indices{^\nu_\mu} x'^\mu
		\end{equation}
		It is conventional to denote the inverse as a lower index prior to
		the one upstairs. This is done to denote that the inverse of a Lorentz transformation
		is also a Lorentz transformation.
		\[
			\Lambda\indices{_\mu^\nu}
			=
			(\Lambda\inv)\indices{^\mu_\nu}
		\]
		\[
			x^\nu  = \Lambda\indices{_\mu^\nu} x^\mu
		\]
		These are just numbers, remember. So, we can just think of these 
		operations as matrix multiplications.

		It is evident that the Lorentz transformations combined
		with its inverse is an identity map
		\begin{align*}
			\Lambda\indices{_\alpha^\beta}
			\Lambda\indices{^\alpha_\gamma}
			&=
			(\Lambda\inv)\indices{^\beta_\alpha}
			\Lambda\indices{^\alpha_\gamma}\\
			&=
			\delta\indices{^\beta_\gamma}
		\end{align*}

		Since $x'\mu$ is a curve, we can get the basis vector of the 
		curve, and the $x^\nu$ curve. 
		\[
			\pdv{}{{x'}^\mu} \equiv \p'_\mu
		\]
		and the $x^\mu$ curve.
		\[
			\pdv{}{x^\mu} \equiv \p_\mu 
		\]

		We can feed this basis to the Minkowski tensor to get the components
		\[
			\eta(\p'_\mu,\p'_\nu) = \eta'_{\mu\nu}
		\]
		and 
		\[
			\eta(\p_\mu,\p_\nu) = \eta_{\mu\nu}
		\]
		However, by the definition of the Minkowski metric, the primed coordinate
		Minkowski metric and unprimed coordinate metric must both be 
		equal to the canonical definition of the Minkowski metric, hence, 
		must be equal to
		\begin{align*}
			\eta'_{\mu\nu} 
			&= \eta_{\mu\nu}\\
			&= \diag(-1,1,1,1)
		\end{align*}
		We can relate them together to see 
		how the basis vectors transform under the change of coordinates,
		recall how a vector transforms:
		\[
			\p'_\mu = \pdv{x^\nu}{x'^\mu}\p_\nu
		\]
		However, this is really nothing but (\ref{eq:lorentz-transformation}), so
		\begin{align*}
			\p'_\mu 
			&= \Lambda\indices{^\nu_\mu}\p_\nu  \\
			&= (\Lambda\inv)\indices{^\mu_\nu}\p_\nu  \\
			&= 
			\Lambda\indices{_\nu^\mu}\p_\nu
		\end{align*}
		We can again feed this basis vector to the tensor $\eta$. Recall that $\eta'_{\mu\nu} = \eta_{\mu\nu}$
		\begin{align*}
		    \eta
			\pqty{
				\p'_\mu,
				\p'_\nu
			}
			&=
			\eta
			\qty(
				\Lambda\indices{_\mu^\alpha}
				\p_\alpha,
				\Lambda\indices{_\nu^\beta}
				\p_\beta,
			)\\
			\eta'_{\mu\nu}
			&=
			\Lambda\indices{_\mu^\alpha}
			\Lambda\indices{_\nu^\beta}
			\eta_{\alpha\beta}\\
			\eta_{\mu\nu}
			&=
			\Lambda\indices{_\mu^\alpha}
			\Lambda\indices{_\nu^\beta}
			\eta_{\alpha\beta}
		\end{align*}

		We have now arrived at a defining property of
		the Lorentz transformation. A Lorentz transformation is
		a linear transformation between inertial coordinate systems
		which preserves the Minkowski metric.

		\begin{definition}{Lorentz Transformation}
			A Lorentz transformation is a linear transformation
			$\Lambda$ which preserves the Minkowski metric, such that
			\begin{equation}
				\label{eq:lorentz-transformation}
				\eta_{\mu\nu}
				=
				\Lambda\indices{_\mu^\alpha}
				\Lambda\indices{_\nu^\beta}
				\eta_{\alpha\beta}.
			\end{equation}
		\end{definition}
		We can write this in a neater way. We represent the equation
		explicitly
		\[
			(\Lambda\inv)\indices{^\alpha_\mu}
			(\Lambda\inv)\indices{^\beta_\nu}
			\eta_{\alpha\beta}
			=
			\eta_{\mu\nu}
		\]
		These are just numbers, so we can treat them just as matrices.
		We can move the $\eta$ term around, and find
		\[
			(\Lambda\inv)\indices{^\alpha_\mu}
			\eta_{\alpha\beta}
			(\Lambda\inv)\indices{^\beta_\nu}
			=
			\eta_{\mu\nu}
		\]
		We can transpose the first matrix to make it match
		the index of $\eta$ (recall the requirement of matrix multiplication)
		\[
			(\Lambda\inv)\indices{^\mu_\alpha}
			\eta_{\alpha\beta}
			(\Lambda\inv)\indices{^\beta_\nu}
			=
			\eta_{\mu\nu}
		\]
		Put symbolically, treating $\Lambda\indices{^\alpha_\mu}$ as 
		simply a matrix, $\Lambda$, and $\eta$ as just another matrix.
		Since the inveres of a Lorentz transformation is itself a
		Lorentz transformation (\ref{eq:lorentz-inverse-transform}),
		we can just replace $\Lambda\inv$ with $\Lambda$.
		Thus, we have the canonical definition of a Lorentz transformation
		\begin{equation}
			\label{eq:lorentz-transformation-symbol}
		\Lambda^\top\eta\Lambda=\eta
		\end{equation}
		or equivalently,
		\[
			\Lambda\eta\Lambda^\top = \eta
		\]
		All Lorentz transformations satisfy this property. The Minkowski
		metric must remain invariant under Lorentz transformations.

		What else can we find from this condition?
		We can try to find the determinant of the matrix $\Lambda$.

		Starting from the Lorentz condition, we have
		\[
			\det(\Lambda^\top \eta \Lambda) = \det(\eta)
		\]
		Since we know $\eta = \diag(-1,1,1,1)$, and we know that
		\[
			\det(ABC) = \det A \det B \det C
		\]
		and the determinant of a transpose is invariant
		\[
			\det(A^\top) = \det(A)
		\]
		we get that
		\[
			\det(\Lambda)^2 \det(\eta) = \det(\eta)
		\]
		Since 
		\[
			\det(\eta) = (-1)(1)(1)(1) = -1\neq 0
		\]
		we can cancel it, and hence we find
		\[
			\det(\Lambda)^2 = 1
		\]
		Which gives us the result
		\begin{equation}
			\label{eq:det-lorentz}
			\det(\Lambda) = \pm 1
		\end{equation}
		This gives us two branches for solutions of Lorentz transformation
		matrices. We call them the following
		\[
			\det(\Lambda) = 
			\begin{cases}
				+1&\quad \to \quad \text{Proper Lorentz transformations}\\
				-1&\quad \to \quad \text{Improper Lorentz transformations}\\
			\end{cases}
		\]

		The former is called a proper Lorentz transformation, because 
		applying two Lorentz proper Lorentz transformations still
		results to a proper transformation. 

		We call the latter an improper transformation because it does 
		not preserve the determinant of the total transformation, since $-1\cdot -1 = 1$,
		an improper transformation applied to another may result to a proper transformation.
		Additionally, if you try to analyze this physically, the improper 
		transformations flip the direction of time. It implies a particle
		moving backwards in time, which is the standard Feynman-St\"uckelberg
		interpretation of antiparticles. However, improper transformations
		don't form a connected group with the identity, or map the physical
		configuration of space to something physically inequivalent.

		Hence, we will only use proper Lorentz transformations
		for physics applications.
	
	\subsubsection{Invariance of $c$ from the metric}
		
		In the last sections, we discussed how the norm of 
		a vector defines its length. A norm is just 
		the square root of the inner product of a vector with itself.
		We will now learn that the norm is in fact invariant, it does
		not depend on the manifold or the coordinates.
		\begin{theorem}{Norm Invariance}
		    For a vector $p$ in a manifold $\mcal M$ equipped
			with a metric $g$, the inner product is invariant 
			under coordinate transformations.
			\begin{equation}
				\label{eq:norm-invariance}
			    g'(p',p') = g(p,p)
			\end{equation}
		\end{theorem}
		

		\begin{proof}
			Consider a manifold equipped with a metric 
			$g$. We have a tangent vector $p$, and its norm
			\[
				\abs{p} = \sqrt{g(p,p)}
			\]
			We can represent $p$ in terms of the coordinate basis
			\[
				p = p^\mu \p_\mu
			\]
			Then, we have
			\[
				g(p,p) = g_{\mu\nu}p^\mu p^nu
			\]
			Under a coordinate transformation,
			\[
				x^\mu\to x'^\mu
			\]
			The components $p^\mu$ and $g_{\mu\nu}$ change such that
			\[
				p'^\mu = 
				\pdv{x'^\mu}{x^\nu} p^\nu 
			\]
			and
			\[
				g'_{\mu\nu} = 
				\pdv{x^\alpha}{x'^\mu} 
				\pdv{x^\beta}{x^\nu} 
				g_{\alpha\beta}
			\]
			When we substitute it into the inner product, we 
			find that 
			\begin{align*}
				g'_{\mu\nu}
				p'^\mu
				p'^\nu
				&=
				\pdv{x^\alpha}{x'^\mu} 
				\pdv{x^\beta}{x^\nu} 
				g_{\alpha\beta}
				\pdv{x'^\mu}{x^\alpha} p^\alpha 
				\pdv{x'^\mu}{x^\beta} p^\beta\\
				&=
				\pdv{x^\alpha}{x'^\mu} 
				\pdv{x'^\mu}{x^\alpha} 
				\pdv{x^\beta}{x^\nu} 
				\pdv{x'^\mu}{x^\beta}
				g_{\alpha\beta}
				p^\alpha 
				p^\beta\\
				&= 
				g_{\alpha\beta}
				p^\alpha
				p^\beta
			\end{align*}
			Therefore, the coordinate transformation did not change 
			the resulting inner product;
			\[
				g'(p',p') = g(p,p)
			\]
			Hence, the norm of a vector is invariant.
		\end{proof}
		
		In relativity, we refer to the norm, or the distance 
		between to points in spacetime, as the space time interval.
		\[
			ds^2 = \eta_{\mu\nu}\dd{x^\mu}\dd{x^\nu}
		\]
		Since we have just shown that the norm is invariant, it then
		folows that the spacetime interval is also invariant. 

		Light is a photon; its wordline is known as a null geodesic, or
		a null curve,
		which means that the spacetime interval of the curve
		along its worldline vanishes.

		For a photon $\gamma$,
		\[
			0 = \eta_{\mu\nu}\dd{x^\mu}\dd{x^\nu}
		\]
		
		In Lorentz coordinates $(ct,x,y,z)$, we will find that this defines 
		the norm of a light curve to be
		\[
			0 = -c^2 \dd{t}^2 + \dd{x}^2 + \dd{y}^2 + \dd{z}^2
		\]
		Invariance implies that
		\[
			c^2 \dd{t}^2 = \dd{x}^2 + \dd{y}^2 + \dd{z}^2
		\]
		If we divide by $\dd{t}^2$, we find
		\[
			c^2 = 
			\pqty{
				\dv{x}{t} 
			}^2
			+
			\pqty{
				\dv{y}{t} 
			}^2
			+
			\pqty{
				\dv{z}{t} 
			}^2
		\]

		If we transform coordinates from $x\to x'$ and $y\to y'$ and so on,
		noting that a transformation in Euclidean coordinates is trivial,
		we find that
		\[
			c^2 = 
			\pqty{
				\dv{x'}{t} 
			}^2
			+
			\pqty{
				\dv{y'}{t} 
			}^2
			+
			\pqty{
				\dv{z'}{t} 
			}^2
		\]
		That is, different coordinate systems give the same value.
		What is this value?
		This is basically just the dot product (normal Euclidean space) of the typical velocity
		\[
			c^2 = \vb{v}\cdot \vb{v} = v^2
		\]
		Which means that we get the speed of a particle travelling 
		at a null curve, which is
		\[
			c = v
		\]
		That is, all inertial observers would measure the speed of a particle
		travelling along a null geodesic, such as a photon, to be the same.
		Thus, the speed of light $c$ is invariant.

		We have already established this result in \S\ref{subsec:relativity-without-light},
		but now we see that the invariance of $c$ is embodied in the invariance
		of the metric.

		While this result emerges naturally from the elementary assumptions
		of relativity, we now see that the invariance of $c$ also follows
		directly from the Minkowski metric.
		
	\subsubsection{Particle velocity bounds}
		\label{subsubsex:particle-velocity-bounds}

		Consider a timelike curve $\gamma$, with tangent vector $u$.
		Let $u$ be parameterized by $t$, or $x^0$, so that its tangent 
		vector is 
		\[
			u = \dv{x^\mu}{t}\p_\mu 
		\]
		The metric then gives
		\[
			g(u,u) =
			-c^2 
			+
			\pqty{
				\dv{x}{t} 
			}^2
			+
			\pqty{
				\dv{y}{t} 
			}^2
			+
			\pqty{
				\dv{z}{t} 
			}^2
		\]

		A timelike curve, which was discussed earlier, has a spacetime interval
		that is negative, i.e.,
		\[
			g(u,u) < 0
		\]
		Therefore 
		\[
			-c^2 + 
			\pqty{
				\dv{x}{t} 
			}^2
			+
			\pqty{
				\dv{y}{t} 
			}^2
			+
			\pqty{
				\dv{z}{t} 
			}^2
			< 0
		\]
		Once again, the spatial terms is just the normal velocity $v$, so we have
		\[
			- c^2 + v^2 < 0
		\]
		Rearranging and taking the square root, we have
		\[
			v < c
		\]
		That is, for a timelike curve, a particle must always have a speed $v<c$.

		Using a similar argument for spacelike curves, $g(u,u)> 0$, we obtain that they must have
		a speed $v > c$. 
		
		We can further prove the following (See \S\ref{}), that 
		for massive particles, i.e, particles which have an invariant
		mass, the curves they travel are always timelike. Similarly,
		massless particles always travel null geodesics. Finally,
		superluminal particles, which are hypothesized but not detected,
		always spacelike.
		Hence, we have the following relationship

		\begin{align*}
			\text{massive particles}\quad &\implies \quad 
			\text{timelike}\quad&\implies v< c\\
			\text{massless particles}\quad &\implies \quad 
			\text{null}\quad&\implies v= c\\
			\text{superluminal particles}\quad &\implies \quad 
			\text{spacelike}\quad&\implies v> c
		\end{align*}
	
		For an observer, timelike curves are worldlines that lie
		\textit{within} the light cone, while null curves are worldlines
		that lie along the light cone, i.e., it is tangent to the light cone. Finally, hypothetical superluminal
		particles, or spacelike curves, have worldlines that lie
		outside the light cone.
		
		\begin{figure}[h]
			\centering
			\includegraphics[width=3in]{images/wordline-lightcone.pdf}
			\caption{Worldlines of timelike, null, and spacelike curves on the light cone}
			\label{fig:wordline-lightcone}
		\end{figure}

		One might think that spacelike curves cannot represent physical
		motion, but they are mathematically allowed within the geometry
		of Minkowski spacetime. The speed of light is not a hard limit
		on all possible curves; rather, it separates the causal
		structures of spacetime. Massive particles must always travel at
		speeds less than $c$, while massless particles travel at exactly
		$c$. A hypothetical superluminal particle, however, would always have
		to travel at speeds greater than $c$.

		For spacelike curves, it becomes difficult to assign
		an absolute distinction between past and future. Since spacelike
		separated events cannot be causally connected, different inertial
		observers can disagree about their temporal ordering. In this
		sense, there is no invariant notion of which event occurs first.
		Additionally, for superluminal particles, we also have difficulty
		assigning the concept of mass. If we take the mathematics literally,
		it would show that superluminal particles (called Tachyons) can either have
		negative mass, or travel in imaginary time.

		Nothing has been empirically observed before. 
		Thus, we mostly think it's not physical.  

		\begin{remark}
			On Lorentz coordinates, we often see a factor of $c$
			in front of the time coordinate,
			\[
				(ct,x,y,z).
			\]
			This is so that the $t$-coordinate has units of length,
			making all four coordinates dimensionally consistent.

			\hspace{1em}It is a common notion in popular science that time is
			the fourth dimension. More precisely, what we see here is
			that spacetime is a 4D manifold equipped with the Minkiwski
			metric, embedded in $\R^4$. It is all spatial.

			\hspace{1em}There is also a convention in which we set $c=1$, so that 
			time has a unit of length, and
			the metric simply becomes.
			\[
				ds^2 = - dt^2 + dx^2 + dy^2 + dz^2
			\]
			This is known as using natural units. It is a unit 
			of convenience, since $c$ is a prefactor that is present
			everywhere. It simplifies
			calculations, and we will adopt
			this convention from time to time.
		\end{remark}

	\subsubsection{Proper Time}	
		\label{subsubsec:proper-time}

		We will now discuss what is arguably one of the 
		most important invariant in general relativity.

		Consider the spacetime interval $\dd{s}$. It has value
		\begin{align*}
			\dd{s^2} &= g_{\alpha\beta}\dd{x^\alpha}\dd{x^\beta}\\
					 &=\eta_{\alpha\beta}\dd{x^\alpha}\dd{x^\beta}
		\end{align*}
		Suppose
		we have a timelike curve $\gamma$ in spacetime, and we wish 
		to find the arclength of this curve from a point $A$ to
		another point $B$\footnote{Arclength is a term we use here, since it's 4D, it's mostly an analogy}.
		We get that.
		\[
			\ell(\gamma) 
			=
			\int_A^B \sqrt{-\dd{s^2}}
		\]
		We put $-\dd{s^2}$ under the radicand since for timelike curves,
		$\dd{s^2}<0$.

		Suppose we wish to know the time traversed by a particle
		to get from point $A$ to $B$. Since $\dd{s}$ is a distance
		(recall that all the components of the Lorentz coordinates
		have units of distance), a natural thing to do to obtain a
		time is to divide a distance by a velocity. The natural
		velocity to use is the invariant speed $c$.
		\[
			\tau = \frac{\lambda}{c} 
		\]
		Thus, we have this \textit{invariant time} of the spacetime
		arclength, which is the time a particle moving 
		in spacetime traverses to get from an event $A$ to an event $B$.
		\[
			\tau \equiv \frac{1}{c} 
			\int_A^B \sqrt{-\dd{s^2}}
		\]
		Suppose that we parametrize $\gamma$ with a parameter $\lambda$,
		since curves require parameters, and a world line is a curve.
		The tangent vector of $\gamma$ is
		\[
			\dot\gamma = \pdv{}{\lambda} 
		\]
		We can then substitute that in the definition of $\tau$, so 
		that we have an even more concrete relation
		\[
			\tau = 
			\frac{1}{c}\int_{\lambda_A}^{\lambda_B} 
			\sqrt{
				-\eta\qty(
				\pdv{}{\lambda},
				\pdv{\lambda}
				)
			}
			\dd{\lambda}
		\]
		But $\pdv{}{\lambda}$ in terms of its components is just
		\[
			\pdv{}{\lambda} 
			=
			\dv{x^\alpha}{\lambda} 
			\pdv{}{x^\alpha} 
		\]
		Hence
		\begin{align*}
			\tau = 
			\frac{1}{c}\int_{\lambda_A}^{\lambda_B} 
			\sqrt{
				-
				\eta_{\alpha\beta}
				\dv{x^\alpha}{\lambda}
				\dv{x^\beta}{\lambda} 
			}
			\dd{\lambda}
		\end{align*}

		We can explicity expand the term inside the radicand
		with the Lorentz coordinates $x^\alpha = (-ct,x,y,z)$,
		to get
		\[
			\tau = 
			\frac{1}{c}\int_{\lambda_A}^{\lambda_B} 
			\sqrt{
				c^2 
				\qty(
					\dv{t}{\lambda} 
				)^2
				-
				\qty(
					\dv{x}{\lambda} 
				)^2
				-
				\qty(
					\dv{y}{\lambda} 
				)^2
				-
				\qty(
					\dv{z}{\lambda} 
				)^2
			}
			\dd{\lambda}
		\]
		We can take the differential of $\tau$, and by the first
		fundamental theorem of calculus, we get
		\[
			\dd{\tau} 
			=
			\frac{1}{c}
			\sqrt{
				c^2 
				\qty(
					\dv{t}{\lambda} 
				)^2
				-
				\qty(
					\dv{x}{\lambda} 
				)^2
				-
				\qty(
					\dv{y}{\lambda} 
				)^2
				-
				\qty(
					\dv{z}{\lambda} 
				)^2
			}
			\dd{\lambda}
		\]
		We can square both sides to get
		\[
			\dd{\tau}^2
			=
			\frac{1}{c^2}
			\pqty{
				c^2 
				\qty(
					\dv{t}{\lambda} 
				)^2
				-
				\qty(
					\dv{x}{\lambda} 
				)^2
				-
				\qty(
					\dv{y}{\lambda} 
				)^2
				-
				\qty(
					\dv{z}{\lambda} 
				)^2
			}
			\dd{\lambda}^2
		\]
		But this is just
		\[
			\dd{\tau^2}
			=
			-
			\frac{1}{c^2}
			\pqty{
				-c^2\dd{t}^2
				+
				\dd{x}^2
				+
				\dd{y}^2
				+
				\dd{z}^2
			}
		\]
		In Lorentz coordinates, $\dd{s}^2$ is defined to be
		\begin{align*}
			\dd{s}^2 
			&= 
			\eta_{\alpha\beta}\dd{x^\alpha}\dd{x^\beta}\\
			&=
				-c^2\dd{t}^2
				+
				\dd{x}^2
				+
				\dd{y}^2
				+
				\dd{z}^2
		\end{align*}
		Hence, the value of $\dd{\tau}^2$ is just
		\[
			\dd{\tau^2}
			=
			-
			\frac{\dd{s^2}}{c^2} 
		\]
		or, if one prefers,
		\[
			\dd{\tau^2}
			=
			\dd{t}^2
			-
			\frac{
				\dd{x}^2
				+
				\dd{y}^2
				+
				\dd{z}^2
			}{c^2} 
		\]
		Notice that this quantity is invariant. 
		It does not change with different observers.
		Switching coordinates will vary $t,x,y,$ and $z$,
		but $\tau$ remains invariant. That is,
		all observers will agree on the value of $\tau$. 

		\paragraph{Comoving observer}

		A comoving observer is an observer that travels along
		with a particle along the same world line. They move with
		the particle with zero displacement. 

		Since they are an observer, they are allowed
		to make measurements of the spacetime around them.
		For a comoving observer, the particle has no relative 
		displacement. To him, the values of the spatial components are
		\[
			\dd{x} = 0,\quad \dd{y}=0,\quad \dd{y} = 0
		\]
		Thus, for a comoving observer, they will measure $\tau$ as simply
		\[
			\dd{\tau} 
			\overset{*}{=}
			\dd{t}
		\]
		The $*$ there signifies the comoving observer's measurement.
		Thus, $\tau$ is precisely the time measured by a clock traveling 
		along a world line. It is the measurement the elapsed time of events 
		that occured on the world line. It is what a comoving observer would measure as 
		the time that has elapsed on their reference frame.

		By definition, you are comoving with yourself. Hence, $\tau$
		is precisely the elapsed time that you can measure between events happening 
		in your own world line.

		Since it is invariant, non-comoving observers will measure different 
		$\dd{t}$ and  $\dd{x}$ and so on, but they will measure the same $\dd{\tau}$.
		They will all agree that an event happened in $\tau$ amount of time from
		their own observations.
		Hence, it is a special invariant quantity.
		We call this measurement the \textit{proper time}.\footnote{The world used in original German 
		for proper time is ``\textit{eigenzeit}''. The word \textit{eigen} means `proper', or `to one's own',
		the same root word for \textit{property}. This is the same root used in the terms \textit{eigenvector} and
		\textit{eigenvalue}, since those quantities are inherent, or owned only, by a matrix itself. That definition of proper 
		has fallen in modern usage, so it's a bit difficult to understad why we call it that today. It's a classical
		terminology. As such, \textit{proper time} is the time experienced by an observer in their own world line.}

		\begin{definition}{Proper Time}
		    The proper time $\tau$ is an invariant quantity
			that an
			observer will measure as the elapsed time in their local world line
			and reference frame.
			It is defined to be
			\begin{equation}
				\label{eq:proper-time}
				\dd{\tau}^2 = -
				\frac{1}{c^2} \eta_{\mu\nu}\dd{x^\mu}\dd{x^\nu} 
			\end{equation}
			External observers can calculate $\dd{\tau}$
			from measurements taken from their respective worldlines, and
			all will agree on its value.
		\end{definition}

		\paragraph{Affine parameter}

		The quantity $\tau$ is an affine parameter. It can be chosen
		as a parameter for the curve $\gamma$, for which the geodesic
		equation takes its standard form, and lienar transformations on
		it preserve the geodesic equation. 
		Any other affine parameter is related to it by a linear transformation.
		For example, if  
		$
			\tau' = 	a \tau + b
		$,
		then $\tau'$ is also affine.
		


		% The parameter $\lambda$
		% describes the curve, while the quantity inside the square root is
		% the contraction of the tangent vector with itself under the metric.
		% Under a change of coordinates, both the components of the tangent
		% vector and the components of the metric change, but $\tau$
		% does not. Thus, the value of $\tau$ is independent of the coordinates
		% used to describe the curve.

		% That is, we can change the parameter from $\lambda$ to be $\tau$

		% Thus, $\tau$ is precisely the time measured by a clock traveling
		% along the worldline. Since it was constructed from the invariant
		% spacetime interval, this elapsed time is itself invariant.
	\subsubsection{Four-velocity, Four-momentum}
		\paragraph{Four-velocity}
		When we defined $\tau$, we used an arbitrary parameter $\lambda$.
		However, since $\tau$ is invariant of the coordinates and is an affine
		parameter, we might as well just use $\tau$
		as the parameter for the curve. We can then define the tangent vector 
		of a curve $\gamma$ to be
		\[
			\dot{\gamma} = \pdv{}{\tau} 
		\]
		Or, when expanded with coordinates $\qty{x^\mu}$, it becomes
		\[
			\pdv{}{\tau} = \dv{x^\mu}{\tau} \pdv{}{x^\mu}   
		\]
		This is a vector. Since it is spatial coordinates differentiated with \textit{time},
		we can therefore think of it as a velocity. Therefore, we have
		the velocity of the world line to be
		\[
			u = \pdv{}{\tau} 
		\]
		We can expand this using the coordinates to be
		\begin{align*}
			u &= \pdv{x^\mu}{\tau} \pdv{}{x^\mu}  \\
			  &= \dot x^\mu \p_\mu
		\end{align*}
		This velocity has four dimensions, which is 
		how each of $(t,x,y,z)$ changes over time. This is
		what we call the \textbf{Four-velocity}.
		\begin{definition}{Four-velocity}
		    The four-velocity of a particle with a worldline 
			$x^\mu(\tau)$ is the tangent vector to that world line
			with respect to proper time,
			\[
				u = \pdv{x^\mu}{\tau} \p_\mu 
			\]
		\end{definition}
			
		To differentiate with the usual velocity, the three-velocity
		\begin{align*}
		    v &= 
			\dv{x^\alpha}{t}\p_\alpha\\
			 &=
			\dv{x^\alpha}{x^0}\p_\alpha
		\end{align*}
		we denote the three-velocity as $v_{(3)}$ and the four-velocity
		as $v_{(4)}$. However, when you see something like $v^a$ (AIN), it is
		usually taken to be the four-velocity. 

		We can take the norm of the four-velocity, and we can show that
		\begin{align*}
			\eta(v, v)
			&=
			\eta_{\mu\nu} v^\mu v^\nu \\
			&= 
			\frac{-c^2 dt^2 + dx^2 + dy^2 + dz^2}{d\tau} \\
			&=
			\frac{ds^2}{d\tau^2} \\
			&=
			\frac{-c^2 \qty(-\dfrac{ds^2}{c^2})}{d\tau^2} \\
			&= -c^2 \frac{\dd{\tau}^2}{\dd{\tau}^2} \\ 
			&= -c^2
		\end{align*}
		The magnitude is thus $\abs{v_{(4)}}= c$. Hence, all particles travel with speed $c$ in space time. 
		We are all moving at the speed of light.

		This is another interpretation of why nothing can go faster than
		light, because all particles are already moving at the speed of light
		in space time. If we are to define $\vb{x} = (x,y,z)$ and define $\dd{s}$ from there,
		we find that
		\[
			ds^2 = \dd{\vb{x}}^2 - c^2\dd{t^2}
		\]
		Photon velocity, or lightlike curve velocity, has norm $\dd{s} = 0$. In spacetime is maximally spatial, which means that
		all of its four-velocity components is spatial, and its temporal component
		is 0, $\dd{t}=0$, which is why the proper time of photons is 0. 
		\begin{align*}
			ds^2 &= \dd{\vb{x}}^2 - c^2\dd{t^2}\\
			0 &= \dd{\vb{x}}^2 - c^2\dd{t^2}\\
			c^2 \dd{t}^2 &= \dd{\vb{x}}^2\\
			c^2 &= \frac{\dd{\vb{x}^2}}{\dd{t}^2} \\
			c &= \dv{\vb{x}}{t} 
		\end{align*}
		A particle at rest 
		has maximal temporal velocity, which means it is traveling through
		space time with 0 spatial velocity, i.e., $\dd{\vb{x}}/\dd{t}=0$. It therefore perceives proper time
		as simply its own time. 
		\begin{align*}
			ds^2 &= \dd{\vb{x}}^2 - c^2\dd{t^2}\\
			ds^2 &= -c^2\dd{t^2}\\
			-\frac{ds^2}{c^2} &= \dd{t^2}\\
			\dd{\tau}^2 &= \dd{t^2}\\
			\dd{\tau} &= \dd{t}
		\end{align*}

		A moving particle has components of spatial
		and temporal velocity, and since the four-velocity is constant, it has 
		to trade off between components of spatial and temporal velocity. Hence, 
		the greater its spatial velocity (the greater the speed), the lower the temporal
		velocity (the slower time flows). This is where the phenomenon of time dilation 
		at high speeds originate. We will discuss this in detail in a bit. 

		We need a way to switch between the 4D spacetime picture,
		and the usual 3D picture we are more familiar to. We defined
		the four-velocity to be
		\[
			u^a = \pdv{}{\tau} 
		\]
		We can expand its components 
		\[
			u^a = u^\alpha \p_\alpha 
		\]
		and define those components to be the Lorentz coodinates.
		\[
			u^a = \pqty{
				c\dv{t}{\tau},
				\dv{x}{\tau},
				\dv{y}{\tau},
				\dv{z}{\tau},
			}
		\]
		Let $\vb{x} = x^i = (x,y,z)$. The $i$ here does not denote AIN, but rather denotes spatial components.
		We can therefore compactify the notation to be
		\[
			u^a = \pqty{
				c \dv{t}{\tau},
				\dv{\vb{x}}{\tau} 
			}
		\]
		We can use chain rule to pull out $\dv{t}{\tau}$ on the spatial components
		\[
			u^a 
			=
			\pqty{
				c \dv{t}{\tau},
				\dv{\vb{x}}{t} \dv{t}{\tau}  
			}
		\]
		And we can pull out the $\dv{t}{\tau}$ on the entire expression.
		\[
			u^a 
			=
			\dv{t}{\tau} 
			\pqty{
				c,
				\dv{\vb{x}}{t} 
			}
		\]
		Notice that on the left hands side is $u^a$, a four velocity, $v_{(4)}$,
		and on the right hand side is the components of the three velocity, $v^i_{(3)}$.
		Hence, to convert between $v_{(4)}$ and $v_{(3)}$, we have
		\[
			v_{(4)}
			=
			\dv{t}{\tau}
			\pqty{
				c,
				v^i_{(3)}
			}
		\]
		While $v_{(4)}$ is invariant in space time, $v_{(3)}$ itself 
		is not invariant. However, since we know the norm of $v_{(4)}$,
		\[
			\eta_{\alpha\beta}u^\alpha u^\beta = -c^2
		\]
		We can determine $\displaystyle\dv{t}{\tau}$.
		That is,
		\begin{align*}
			\abs{v_{(4)}} 
			&= 
			-
			\pqty{
				c \dv{t}{\tau} 
			}^2
			+
			\pqty{
				\dv{\vb{x}}{\tau} 
			}^2\\
			-c^2
			&=
			-
			\pqty{
				c \dv{t}{\tau} 
			}^2
			+
			\pqty{
				\dv{\vb{x}}{t} 
			}^2
			\pqty{
				\dv{t}{\tau} 
			}^2
			\\
			&=
			\pqty{
				\dv{t}{\tau} 
			}^2
			\pqty{
				-c^2
			+
				\pqty{
					\dv{\vb{x}}{t} 
				}^2
			}
			\\
			-c^2
			&=
			\pqty{
				\dv{t}{\tau} 
			}^2
			\pqty{
				-c^2
			+
				v^2_{(3)}	
			}\\
			\pqty{
				\dv{t}{\tau} 
			}^2
			&=
			\frac{
				c^2
			}{
				\pqty{
					c^2
				-
					v^2_{(3)}	
				}
			}
			\\
			\dv{t}{\tau} 
			&=
			\sqrt{
				\frac{
					c^2
				}{
					{
						c^2
					-
						v^2_{(3)}	
					}
				}
			}
			\\
			\dv{t}{\tau} 
			&=
			\sqrt{
				\frac{
					1
				}{
					{
						1
						-
						\frac{
							v^2_{(3)}	
						}{c^2} 
					}
				}
			}
			\\
			\dv{t}{\tau} 
			&=
			\frac{
					1
				}{
					\sqrt{
						1
						-
						\frac{
							v^2_{(3)}	
						}{c^2} 
					}
				}
		\end{align*}
		This is just the $\gamma$ factor in Lorentz transformation, or the Lorentz factor.
		It is a function of $v_{(3)}$. Note that $v^2_{(3)}$ is the squared
		norm of $v_{(3)}$, while $v^i_{(3)}$ refers to the components itself.

		Hence, we have 
		arrived with the value of the rate of change of time 
		with respect to proper time:
		\begin{equation}
			\label{eq:dtdtau-gamma}
			\dv{t}{\tau} = \gamma(v_{(3)}) 
		\end{equation}

		This is where the notion that the rate of the passage time of an object in high speed,
		with respect to your proper time (which you measure locally) slows down. As a 
		particle goes faster, $\dd{t}/\dd{\tau}$ becomes lower and lower, until it is $0$ at
		$c$, which means time has completely stopped. 

		This is made clearer by explicitly showing their proportionaly relation.
		\[
			\dd{\tau} = \frac{1}{\gamma}\dd{t} 
		\]

		From this, we also learn how to convert from three-velocity to four-velocity. 
		\begin{equation}
			\label{eq:3vel-to-4vel}
			v_{(4)} = \gamma(v_{(3)}) \pqty{ c , v^i_{(3)} }
		\end{equation}

		Explicity, if $v_x = \dv{x}{t}$, and $v_y$ and $v_z$ are 
		all three-velocities, then
		\[
			v_{(4)} = 
			\frac{1}{
				\sqrt{
					1 - 
					\dfrac{
						v_x^2 + v_y^2 + v_z^2  
					}{c^2} 
				}
			}
			\bigl(
				c, v_x, v_y, v_z
			\bigr)
		\]
		Thus, this formulation bridges the gap between $\R^3$ and Minkowski 
		spacetime. 
		That is, if we know the three-velocity, we can immediately know
		the four-velocity. This will work for any three-velocity $v_{(3)} < c$.
		If the three-velocity is larger than $c$, the conversion will take
		an imaginary value. While it can be studied, it has been done before,
		it does not have any physical meaning. Thus,
		it is deemed to be non-physical. Any object travelling with speed below $c$ will always travel 
		below it, and any object traveling with speed above $c$ can only travel above it.

		\begin{remark}
			$v_{(3)}$ is not a vector in the spacetime manifold, it does not
			deserve an arrow. Only $v_{(4)}$ is a true vector in spacetime. 
		\end{remark}
		
		\paragraph{Four-velocity of null curves}
			Null curves have no proper time. Its proper time is 0. However,
			light still has four velocity, why is that?

			Four velocity is defined to be 
			\[
				u = \dv{x^\alpha}{\tau}\p_\alpha 
			\]
			However, since for null curves, $d\tau=0$ everywhere, then
			\[
				u_\gamma = \frac{\dd{x^\alpha}}{0} \implies \text{undefined} 
			\]
			So there is no four-velocity for null curves in the proper sense.

			Instead, we parameterize the tangent vector of a null curve
			by an affine parameter $\lambda$, which is \textit{not} proper time.
			\[
				k^a = \dv{x^\alpha}{\lambda}\p_\alpha 
			\]
			This is called the wave vector, or null tangent/null vector. 
			Its defining property is:
			\[
				\eta_{ab} k^a k^b = 0 
			\]
			This is the key difference between a massive and 
			massless particle. $k^a$ plays the same strutural
			role as $u^a$.

			It is tangent to the worldline, it satisfies the geodesic equation, i.e.,
			\[
				\dv{k^\sigma}{\lambda} + \Gamma^{\sigma}_{\mu\nu} k^\mu k^\nu = 0 
			\]
			It transforms as a four-vector under Lorentz transformations,
			and its components give frequency and direction.
			\[
				k^\alpha = (\omega/c, \vb{k})
			\]
			It looks like a four-velocity, and in casual usage people refer
			to it as such, but strictly speaking, it's a null vector.
			Its normalization condition is different and proper time plays no role.
			Since the speed of a lightlike curve is constant, the affine parameter
			has no physical scale. It is arbitrary. What is physical is the direction
			of the null vector, and the curve it traces, not its magnitude. The coordinate
			speed of light is always $c$.

			The speed of light comes from the \textit{null condition}.
			\[
				\eta_{\alpha\beta} k^\alpha k^\beta = 0
			\]
			Writing the components as $k^\alpha = (k^0, \vb{k})$, we have
			\[
				-(k^0)^2 + \norm{\vb{k}}^2 = 0 \implies \norm{\vb{k}} = k^0
			\]
			So the coordinate speed is:
			\[
				\frac{\norm{\vb{k}}}{k^0} = 1
			\]
			This is for units of $c=1$
			The ratio is scale-invariant. i.e., $k^a \to \frac{1}{p}k^a, p\in\R$ does not
			change it. So the speed $c$ is a consequence of the null condition. For massive
			particles, proper time \textit{fixes} the scale of $u^a$, because its squared norm 
			is always $-c^2$. Thus, it has a physical magnitude. For a photon, analogous
			normalization is 0, which fixes nothing. Hence, a photon has 
			no preferred scale.

			Light, or any other massless particles, have a world line that is a \textit{null geodesic}.
			That is, the path along which light rays propagate in 
			spacetime follow a null geodesic. It is null because $d\tau=0$ along it, no proper
			time elapses. It is a geodesic because it is force-free, following the straightest
			available path.
			\begin{remark}
				\textbf{Notation}\hspace{1em}
				Indicating a four and three vector is cumbersome.
				So we will adopt the following notation.
				When you see $u^a$ or $v^a$, it is the four-velocity,
				similarly for $p^a$ being momentum.
				We reserve bold-face fonts to represented a three-vector,
				and the index $i$ to give the spatial components of 
				a four-vector (which is kinda equivalent to the three vector).
				Hence, when you see a bold face $\vb{v}$,
				it's a three-velocity. When you see $v^i$, it's a either components of the three velocity, or the
				spatial components of the four-velocity.
				When you see $x^i$, it's the spatial component of 
				a spactime position, and $\vb{x}$, is the three-position. 
				When you see $\vb{p}$, it's the three-momentum, and $p^i$ is either components 
				of the three-momentum, or spatial components of the four-momentum.

				We will specify what $i$ denotes depending on what we work with. Just keep this in mind.
			\end{remark}
		\paragraph{Four-momentum}
			
			The four-momentum, or just momentum\footnote{Sir goes on a rant here that it really should not be called the four-momentum, it's just the momentum. The 3-momentum is just the spatial components of the momentum}
			is the product of a quantity called the 
			invariant mass\footnote{This is just normal inertial mass which
			is invariant for observers. There is no such thing called 
			relativistic mass} and the four-velocity.
			
			\begin{definition}{Four-momentum}
			    For a particle with mass $m$ and four-velocity $u^a$, the 
				four-momentum is
				\[
					p^a = mu^a
				\]
			\end{definition}
			
			
			If you solve for the squared-norm of the four-momentum, we have
			\[
				\eta_{ab} p^a p^b = -m^2c^2
			\]
			The first component of the four-momentum, $p^0$ is:
			\begin{align*}
			    p^0
				&=
				mu^0\\
				&= 
				m\gamma c
			\end{align*}
			To make sense of what that is, 
			let's try multiplying it by $c$.
			\begin{align*}
				cp^0 
				&= m\gamma c^2 \\
				&= m \frac{1}{\sqrt{1-\frac{v^2}{c^2}}} c^2
			\end{align*}
			We can do an expansion to show that 
			\begin{align*}
				m\gamma c^2 &= 
				mc^2 \qty(1 - \frac{1}{2}\frac{v^2}{c^2}+\cdots )\\
							&= 
				mc^2 + \frac{1}{2}mv^2 + \cdots 
			\end{align*}
			Here you will notice that the 
			first term is the well known \textit{rest-mass} energy,
			and the second term is the nonrelativistic kinetic energy,
			and the rest are higher order corrections, and hence we find
			that $m\gamma c^2$ is just the total relativistic energy of a particle.
			\[
				m\gamma c^2 = E
			\]
			Hence, the first component of the four-momentum is
			\[
				p^0 = \frac{E}{c} 
			\]
			For the spatial components of the momentum, $p^i$, $i=1,2,3$, we have
			\begin{align*}
			    p^\mu
				&=
				m\gamma v^i\\
				&=
				\gamma \vb{p}\\
				&=
				\vb{p} \pqty{1 + \frac{1}{2} \frac{v^2}{c^2} + \cdots}
			\end{align*}
			The spatial components of the four-momentum are what is colloquially 
			known as the \textit{relativistic momentum}.

			Thus, the components of the four-momentum is
			\[
				p^\alpha = 
				\pqty{
					\gamma mc, \gamma m v^i
				}
				=
				\pqty{
					\frac{E}{c},
					p^i
				}
			\]

			\paragraph{Conservation}

			Four-momentum is conserved in any interaction. That is,
			\[
				\sum p^a = \sum p^b
			\]
			Where $p^a$ is the momentum of the system before the interaction,
			and $p^b$ is the momentum of the system after.

			This single statement implies both the conservation of energy,
			and the conservation of 3-momentum.
			
			% $p^a$ transforms as a four-vector under Lorentz transformations, just
			% like any other. You can therefore read off $E$ and $\vb{p}$ in any new frame. More on this later.
			% Free particles parallel transport momentum along geodesics.

			The norm is invariant, i.e., frame independent. Any 
			observer measuring the norm of the four-momentum of a particle
			will agree on the norm of the momentum. 
			\[
				\eta_{ab} p^a p^b = -m^2 c^2
			\]
			This is the mass-shell condition (see Remark at the end of the subsection).
			It's the same in every frame. In fact, this is
			how the \textit{invariant mass} is defined in relativity. The momentum is taken as fundamental, and the mass is 
			defined from the invariant four momentum.
			\[
				m = \sqrt{-\frac{1}{c^2} \eta_{ab} p^a p^b}
			\]
			
	\subsubsection{Energy}
	Since the norm is $\norm{p} = -mc^2$, we can show that
			\[
				\eta_{\alpha\beta} p^\alpha p^\beta = - m^2c^2
			\]
			Expanding this out, we have
			\[
				-\pqty{\frac{E}{c^2}} + (p^i)^2 = -m^2c^2
			\]
			Solving for $E$, we arive with the famous formula
			every undergraduate can recite by heart.\footnote{We kinda just contract $p^i\to p$ here to denote the spatial components of the momentum, to get the well known form.}
			\begin{equation}
				\label{eq:energy-momentum relation}
				E^2 = p^2 c^2 + m^2 c^4
			\end{equation}
			With a nonmoving particle, $p=0$. Hence, you get
			the celebrated equation that Einstein is known for the most: The mass-energy equivalence.\footnote{Finally I got to write $E=mc^2$ in an offical capacity. Hell yeah.}
			\begin{equation}
				\label{eq:mass-energy-equivalence}
			    E = mc^2
			\end{equation}

			\paragraph{Measuring energy}
			
			How does one measure the energy of a particle?

			Suppose we have two world lines, $\gamma_0$, which is the
			observer, and $\gamma_1$ the particle being observed.
			How does one measure the energy of $\gamma_1$?.

			Doing what we have learned previously,
			it's kinda cumbersome, is it not?

			Naively, you'd have to
			\begin{enumerate}
			    \item Set up a Lorentz frame adapted to the observer $\gamma_0$.
				\item Find the basis vectors orthogonal to $t$.
				\item Express the particles four velocity $u^a$ in that frame's coodinates.
				\item Determine the mass of the particle.
				\item Combine them to get the energy.
			\end{enumerate}

			This works, but it's clunky. It forces you to pick coordinates
			and do book keeping which hides the geometry.
			
			\begin{figure}[h]
				\centering
				\includegraphics[width=4in, height = 3in]{images/two_worldlines.pdf}
				\caption{Two worldlines, $\gamma_0$ and $\gamma_1$. The observer $\gamma_0$
				measures $\gamma_1$ from its own frame}
				\label{fig:two-worldlines.}
			\end{figure}

			However, knowing the energy-momentum relation, this process
			is simplified. We just have to get the 0-th component of the four momentum?
			How do we do this? We just feed it the $0$-th basis vector, the time basis.
			\[
				p^a(\p_t) = \frac{E}{c} 
			\]
			How do you get the 0th basis vector? Well, it's the time basis vector,
			so the observer, $\gamma_1$ also has it. All he needs to do
			is project the particles four-momentum with his four velocity.
			\[
				E = -\eta_{ab} p^a u^b
			\]
			Where $p^a = mu_{0}^a$, the particle's four momentum, and $u^b$ is
			the observer's four velocity, no frame-building required. 

			Why does this work?

			You are always comoving with yourself, your four velocity
			$u^b$ is always tangent to your world line by definition.	
			The observer's rest frame is defined by the tangent velocity $u^b$.
			Therefore, your four-velocity's spatial coordinates 
			is always zero. Your world line, in your perspective,
			is maximally temporal, and therefore the only component
			in your four-velocity is the time basis vector, $u^b / c$,
			and the spatial basis vectors are any three vectors orthogonal 
			to it. That is, the observer's tangent vector is, by definition,
			the time basis vector.
			\[
				\eta_{ab} p^a u^b = -\frac{E}{c}
			\]
			Projecting $p^a$ in the observer's time direction therefore
			extracts exactly the enregy that the observer assigns to the particle.

			To put it more rigorously, in the observer's rest frame, the 
			four velocity is
			\[
				u^\beta = (c, 0, 0 ,0)
			\]
			So, 
			\begin{align*}
				\eta_{ab} p^a p^b 
				&= \eta_{00} p^0 u^0\\
				&= -1 \qty(\frac{E}{c}) (c)\\
				&= - E
			\end{align*}
			Hence:
			\[
				E = -\eta_{ab}p^a p^b
			\]
			This is the answer to the question ``What will an observer measure?".
			Given an observer, with four-velocity $u^b$, the energy they assign
			to the particle with four momentum $p^a$ is the scalar $E = -\eta_{ab} p^a u^b$.

			The momentum $p^a$ transforms as a four-vector under Lorentz transformations, $p'^c = \Lambda^c_a p^a$.
			As such, you can transform the momeuntum in another frame and read off both $E$ and the 
			spatial momentum $\vb{p}$ in the new frame, by again just operating it on your tangent vector.

			The deeper point, is that energy is not an intrinsic property of the particle — it's a relationship between 
			the particle and the observer. Different observers, different 
			different measured energies. You can know what others are doing relative to you, 
			but there is no frame-independent ``energy of the particle" independent of who measures it.	
		
			In other words,
			You can know what others are doing, but you can't know what you yourself 
			are doing.
			\begin{question}{Does this only work for constant velocities?}
				The answer is no, it does not. Nowhere in the formulation
				did we assume a constant velocity $u^a$, so the formulation
				in general can work
				for accelerated observers. SR in general, works for accelerating
				observers, there is just more calculus involved.

				\hspace{1em}
				The only restrictions we impose is that we are in Minkowski 
				spacetime, and that all particles are Lorentzian observers,
				where moving particles relative to observers are described
				by Lorentzian transformations.

				\hspace{1em}
				No assumption of constant $u^a$ was ever made.
			\end{question}
			
		
			\begin{remark}
				In particle physics, particles which satisfy the energy-momentum relation are 
				called on-shell or on-mass shell, i.e.,
				\[
					p^\alpha p_\alpha = -m^2 c^2
				\]
				Or in 3-vector form, 
				\[
					E^2 = p^2 c^2 + m^2 c^4
				\]

				The relation $E^2 - p^2 c^2 = m^2 c^4$ defines a hyperboloid
				in energy-momentum space, a \textit{shell} at fixed mass $m$.
				A particle whose four-momentum
				\[
					p^\mu = (E/c, \vb{p})
				\]
				lies on this surface is on-shell. One that doesn't is called 
				off-shell. On-shell particles are real, asymptotic particles.
				An example is a photon or an electron. For a photon, the 
				mass-shell is the lightcone, $p^\mu p_\mu = 0$. It also means 
				that four-momentum is conserved. On-shell is a declaration of 
				timelike curves.
				
				\hspace{1em}Particles that are off-shell violate the energy-momentum
				relation. Examples include virtual particles such as virtual
				photons in QED, and internal propagators.
			\end{remark}

		
